\section*{Bài tập 03 --- Cây quyết định}

Dùng logarithm nhị phân trong entropy: \(\log\) hiểu là \(\log_2\).
Ký hiệu \(H(S)\) là entropy của tập nhãn của \(S\) (theo các lớp xuất hiện trong \(S\)).

%--------------------------------------------------------------------
\subsection*{Bài 1. Dự đoán mùa trong năm}

\textbf{Dữ liệu huấn luyện (12 dòng đầu) và hai mẫu cần dự báo X, Y.}

\begin{center}
{\small
\begin{tabular}{@{}ccccc@{}}\toprule
Mẫu & Thời tiết & Lá cây & Nhiệt độ & Mùa\\\midrule
1 & Nắng & Vàng & Trung bình & Thu\\
2 & Tuyết & Xanh & Thấp & Đông\\
3 & Tuyết & Vàng & Thấp & Đông\\
4 & Mưa & Vàng & Trung bình & Thu\\
5 & Mưa & Rụng & Thấp & Đông\\
6 & Tuyết & Rụng & Thấp & Đông\\
7 & Nắng & Rụng & Thấp & Đông\\
8 & Nắng & Xanh & Trung bình & Xuân\\
9 & Nắng & Xanh & Cao & Hè\\
10 & Mưa & Rụng & Trung bình & Thu\\
11 & Mưa & Xanh & Cao & Hè\\
12 & Mưa & Xanh & Trung bình & Xuân\\
X & Mưa & Vàng & Thấp & ?\\
Y & Tuyết & Rụng & Trung bình & ?\\
\bottomrule
\end{tabular}
}
\end{center}

\subsubsection*{(a) Thuật toán ID3 (entropy)}

\textbf{Bước 1 --- Lần chia đầu (trên cả 12 mẫu).}

Tần số nhãn: Thu (3), Đông (5), Xuân (2), Hè (2).
\[
H(S)= -\tfrac{3}{12}\log\tfrac{3}{12}-\tfrac{5}{12}\log\tfrac{5}{12}
      -2\cdot\tfrac{2}{12}\log\tfrac{2}{12}
      \approx 1{,}8879.
\]

Với mỗi thuộc tính, tính entropy có trọng số sau khi tách theo \emph{tất cả} giá trị của thuộc tính đó, rồi lấy \(\text{Gain}=H(S)-H_{\text{sau tách}}\).

\emph{Thời tiết.} Các nhánh Nắng (4 mẫu), Tuyết (3 mẫu), Mưa (5 mẫu) cho entropy từng nhánh lần lượt khoảng \(2{,}0000\); \(0\); \(1{,}9219\). Trung bình có trọng số
\[
H_{\text{thời tiết}}\approx 1{,}4675,\qquad \text{Gain}\approx 0{,}4204.
\]

\emph{Lá cây.} Nhánh Vàng (3 mẫu: Thu, Đông, Thu), Xanh (5 mẫu), Rụng (4 mẫu) cho
\[
H_{\text{lá}}\approx 1{,}1341,\qquad \text{Gain}\approx 0{,}7538.
\]

\emph{Nhiệt độ.} Nhánh Trung bình (5), Thấp (5), Cao (2): nhánh Thấp và Cao thuần nhất (entropy \(0\)); nhánh Trung bình có \(3\) Thu và \(2\) Xuân nên entropy khoảng \(0{,}9710\).
Vậy
\[
H_{\text{nhiệt độ}}
=\tfrac{5}{12}\,H_{\text{(Trung bình)}}+\tfrac{5}{12}\cdot 0+\tfrac{2}{12}\cdot 0
\approx 0{,}4046,\qquad \text{Gain}\approx 1{,}4834.
\]

Thuộc tính có gain lớn nhất là \textbf{Nhiệt độ} --- chọn làm nút gốc.

\textbf{Bước 2 --- Các lá trực tiếp từ gốc.}
\begin{itemize}\setlength\itemsep{0pt}
  \item Nhiệt độ \(=\) \textbf{Cao}: hai mẫu 9 và 11 đều là Hè \(\Rightarrow\) \textbf{Hè}.
  \item Nhiệt độ \(=\) \textbf{Thấp}: toàn nhãn Đông \(\Rightarrow\) \textbf{Đông}.
  \item Nhiệt độ \(=\) \textbf{Trung bình}: còn 5 mẫu (1,4,8,10,12) với nhãn Thu--Thu--Xuân--Thu--Xuân.
\end{itemize}

\textbf{Bước 3 --- Lần chia trên nhánh «Trung bình».}

Gọi \(S'=\) tập \(5\) mẫu trên. \(H(S')=-\tfrac{3}{5}\log\tfrac{3}{5}-\tfrac{2}{5}\log\tfrac{2}{5}\approx 0{,}9710\).

\emph{Thuộc tính «Thời tiết».} Trong \(S'\) không có mẫu nào có thời tiết Tuyết. Nhánh Nắng gồm (1, Thu) và (8, Xuân). Nhánh Mưa gồm (4, Thu), (10, Thu), (12, Xuân). Entropy có trọng số \(\approx 0{,}9510\); Gain \(\approx 0{,}020\).

\emph{Thuộc tính «Lá cây»:} nhánh «Vàng» gồm (1Thu, 4Thu); «Xanh» gồm (8Xuân,12Xuân); «Rụng» gồm (10Thu) --- đều thuần nhất, entropy từng nhánh là \(0\). Vậy \(H_{\text{sau tách}}=0\), \(\text{Gain}\approx H(S')\approx 0{,}9710\).

Chọn \textbf{Lá cây} trong nhánh này.

\textbf{Cây thu được (mô tả bằng luật if--then).}
\begin{enumerate}\setlength\itemsep{0pt}
  \item Nếu Nhiệt độ \(=\) Cao \(\Rightarrow\) Hè.
  \item Nếu Nhiệt độ \(=\) Thấp \(\Rightarrow\) Đông.
  \item Nếu Nhiệt độ \(=\) Trung bình và Lá cây \(=\) Xanh \(\Rightarrow\) Xuân.
  \item Nếu Nhiệt độ \(=\) Trung bình và (Lá cây \(=\) Vàng hoặc Rụng) \(\Rightarrow\) Thu.
\end{enumerate}

\subsubsection*{(b) Dự báo X và Y}

\begin{itemize}\setlength\itemsep{0pt}
  \item X: (Mưa, Vàng, Thấp) --- theo luật (2) \(\Rightarrow\) \textbf{Đông}.
  \item Y: (Tuyết, Rụng, Trung bình) --- theo luật (4) \(\Rightarrow\) \textbf{Thu}.
\end{itemize}

\subsubsection*{(c) Gini và lỗi phân loại}

Với cùng cách tách đa nhánh như trên, trên bộ 12 mẫu này cả ba tiêu chí (entropy, Gini, lỗi phân loại) đều chọn cùng thứ tự thuộc tính: \textbf{Nhiệt độ} ở gốc, rồi \textbf{Lá cây} trên nhánh «Trung bình»; dự báo X, Y vẫn là \textbf{Đông} và \textbf{Thu}.

%--------------------------------------------------------------------
\subsection*{Bài 2. Độc / không độc của lá (10 mẫu học)}

\textbf{Dữ liệu (Huấn luyện: 1--10; kiểm tra: 11, 12.)}

\begin{center}
{\footnotesize
\begin{tabular}{@{}cccccc@{}}\toprule
Mẫu & Hình dạng & Dạng & Màu & Kích thước & Độc\\\midrule
1 & Tròn & Đơn & Đỏ & Nhỏ & Không\\
2 & Tròn & Kép & Đỏ & Nhỏ & Có\\
3 & Tròn & Đơn & Xanh & To & Không\\
4 & Dài & Đơn & Xanh & Nhỏ & Không\\
5 & Dài & Kép & Vàng & Nhỏ & Có\\
6 & Dài & Đơn & Xanh & Vừa & Có\\
7 & Tròn & Kép & Vàng & To & Không\\
8 & Tròn & Kép & Vàng & Vừa & Có\\
9 & Tròn & Đơn & Đỏ & Vừa & Có\\
10 & Dài & Đơn & Đỏ & Vừa & Có\\
11 & Dài & Kép & Đỏ & Vừa & ?\\
12 & Tròn & Đơn & Xanh & Nhỏ & ?\\
\bottomrule
\end{tabular}
}
\end{center}

Trên \(10\) mẫu, Không và Có có tần \(4\) và \(6\). Entropy gốc
\[
H(S)=-\tfrac{4}{10}\log\tfrac{4}{10}-\tfrac{6}{10}\log\tfrac{6}{10}\approx 0{,}9710.
\]

\subsubsection*{(đ) Theo chỉ số Gini --- lần đầu}

Gini của tập: \(i(S)=1-\bigl(\tfrac{4}{10}\bigr)^2-\bigl(\tfrac{6}{10}\bigr)^2=0{,}48\).

Tách theo \textbf{Kích thước}: nhánh ``Nhỏ'' gồm bốn mẫu 1--2--4--5 với nhãn lần lượt Không, Có, Không, Có; nhánh ``To'' một mẫu Không; nhánh ``Vừa'' ba mẫu đều Có.
Impurity Gini từng nhánh: ``Nhỏ'' là \(1-(\tfrac{2}{4})^2-(\tfrac{2}{4})^2=\tfrac{1}{2}\); ``To'' và ``Vừa'' bằng \(0\) (thuần nhất). Gini có trọng số:
\[
\textstyle \tfrac{4}{10}\cdot 0{,}5+\tfrac{1}{10}\cdot 0+\tfrac{3}{10}\cdot 0=0{,}20,\qquad \Delta\approx 0{,}28
\]
lớn nhất so với các thuộc tính còn lại (xem bảng đối chiếu khi tính tay). Vậy chọn \textbf{Kích thước} ở gốc.

\subsubsection*{Nhánh «Kích thước \(=\) Nhỏ»}

Bốn mẫu còn lại có nhãn Không--Có--Không--Có (entropy \(1\)). Tách theo \textbf{Dạng}: «Đơn» (2 Không), «Kép» (2 Có) --- thuần nhất. Đây là lần chia thứ hai.

Nhánh «To» và «Vừa» đã thuần «Không» và «Có».

\subsubsection*{(e) Dự báo mẫu 11 và 12}

Theo cây tìm được:
\begin{itemize}\setlength\itemsep{0pt}
  \item Mẫu 11 (Dài, Kép, Đỏ, Vừa): Kích thước ``Vừa'' \(\Rightarrow\) \textbf{Có}.
  \item Mẫu 12 (Tròn, Đơn, Xanh, Nhỏ): ``Nhỏ'' rồi ``Đơn'' \(\Rightarrow\) \textbf{Không}.
\end{itemize}

\subsubsection*{(f) Entropy và lỗi phân loại}

Ở gốc, entropy có trọng số sau khi tách theo Kích thước là \(0{,}4000\) (gain \(0{,}5709\)), lớn hơn các thuộc tính khác; lỗi phân loại có trọng số cũng đạt cực tiểu \(0{,}20\) với cùng thuộc tính \textbf{Kích thước}. Bước kế trên nhánh «Nhỏ» vẫn chọn \textbf{Dạng}. Kết quả dự báo mẫu 11, 12 trùng phần (đ)--(e): \textbf{Có} và \textbf{Không}.

%--------------------------------------------------------------------
\subsection*{Bài 3. Quyết định cho vay (10 mẫu học)}

\textbf{Dữ liệu.}

\begin{center}
{\footnotesize
\begin{tabular}{@{}cccccc@{}}\toprule
Mẫu & Phái & Công việc & Học vấn & Độ tuổi & Quyết định\\\midrule
1 & Nam & LĐ Chân tay & Cao đẳng & Trung niên & Không\\
2 & Nữ & LĐ Trí óc & Đại học & Trung niên & Có\\
3 & Nữ & LĐ Chân tay & Phổ thông & Già & Có\\
4 & Nam & LĐ Trí óc & Cao đẳng & Trung niên & Có\\
5 & Nam & LĐ Chân tay & Phổ thông & Thanh niên & Không\\
6 & Nam & LĐ Trí óc & Đại học & Già & Có\\
7 & Nam & LĐ Chân tay & Cao đẳng & Già & Có\\
8 & Nữ & LĐ Chân tay & Phổ thông & Trung niên & Không\\
9 & Nam & LĐ Trí óc & Đại học & Thanh niên & Không\\
10 & Nữ & LĐ Chân tay & Cao đẳng & Già & Có\\
A & Nữ & LĐ Trí óc & Cao đẳng & Già & ?\\
B & Nam & LĐ Chân tay & Phổ thông & Trung niên & ?\\
\bottomrule
\end{tabular}
}
\end{center}

\subsubsection*{(a) Lỗi phân loại}

Tại gốc, lỗi gốc \(=1-\max(\tfrac{4}{10},\tfrac{6}{10})=0{,}40\).
Tách theo \textbf{Độ tuổi}: Già (4 mẫu toàn «Có»), Thanh niên (2 mẫu toàn «Không»), Trung niên (4 mẫu: hai Không và hai Có---lỗi impurity của nút bằng \(1-\tfrac{2}{4}=0{,}5\)). Lỗi có trọng số
\[
\tfrac{4}{10}\cdot 0+\tfrac{2}{10}\cdot 0+\tfrac{4}{10}\cdot 0{,}5=0{,}20,
\]
giảm nhiều nhất so với Phái, Công việc (lỗi có trọng số \(0{,}40\) tại nút đầu) và Học vấn (\(0{,}30\)). Chọn \textbf{Độ tuổi} gốc.

Trên nhánh Trung niên, tách tiếp theo \textbf{Công việc}: hai nhánh đều thuần (Chân tay $\to$ Không; Trí óc $\to$ Có).

Luật (tóm tắt):
\begin{enumerate}\setlength\itemsep{0pt}
  \item Độ tuổi \(=\) Già \(\Rightarrow\) Có.
  \item Độ tuổi \(=\) Thanh niên \(\Rightarrow\) Không.
  \item Độ tuổi \(=\) Trung niên và công việc thuộc nhóm trí óc \(\Rightarrow\) Có.
  \item Độ tuổi \(=\) Trung niên và công việc chân tay \(\Rightarrow\) Không.
\end{enumerate}

\subsubsection*{(b) Dự báo A và B}

\begin{itemize}\setlength\itemsep{0pt}
  \item A (Nữ, Trí óc, Cao đẳng, Già): theo luật (1) \(\Rightarrow\) \textbf{Có}.
  \item B (Nam, Chân tay, Phổ thông, Trung niên): luật (4) \(\Rightarrow\) \textbf{Không}.
\end{itemize}

\subsubsection*{(c) Gini và entropy}

Entropy gốc cũng bằng \(0{,}9710\) như phần 2 (nhãn 4/10 và 6/10 là cặp không--có như nhau).

Tại gốc, entropy có trọng số sau khi tách theo Độ tuổi là \(0{,}4000\) (gain \(0{,}5709\)), lớn nhất; Gini có trọng số \(0{,}2000\) (độ giảm \(0{,}2800\)) cũng đạt cực đại. Bước kế trên «Trung niên» vẫn chọn \textbf{Công việc}. Dự báo A, B giữ nguyên \textbf{Có} và \textbf{Không}.
